\input{../tex-inputs/latex-korrekturansicht-vorspann}

\section[Arthur Schnitzler an Gustav Schwarzkopf, 13. 1. 1925]{L04356 Arthur Schnitzler an Gustav Schwarzkopf, 13. 1. 1925}
\nopagebreak\mylabel{L04356v}
\rehead{ }\normalsize\beginnumbering\briefempfaengerindex{Schwarzkopf, Gustav@\textsc{Schwarzkopf, Gustav}!zzzSchnitzler, Arthur@\emph{von Arthur Schnitzler}!1925-01-131@{13. 1. 1925}|(be}
\toendnotes[C]{\smallbreak\pagebreak[2]}
\correspDesc{Versand  durch Arthur Schnitzler am 13. 1. 1925 in Freiburg im Breisgau
\newline{}Übermittlung  am 13. 1. 1925 in Freiburg im Breisgau
\newline{}Erhalt  durch Gustav Schwarzkopf im Zeitraum [14. 1. 1925 – 18. 1. 1925?] in Wien}\toendnotes[C]{\smallbreak}
\Standort{DLA, A:Schnitzler, HS.1985.1.1897.}
\physDesc{Bildpostkarte, 335 Zeichen
\newline{}Handschrift: Bleistift, lateinische Kurrent
\newline{}Versand: Stempel: »\nobreak{}\oindex{Freiburg im Breisgau@\textbf{Freiburg im Breisgau}, \emph{Hauptstadt}|pwk}Freiburg (Breisgau), 13. 1. 25, 2–3N\nobreak{}«.  }\toendnotes[C]{\smallbreak}\pstart{}{\pb}Hn Gustav Schwarzkopf\pend{}\pstart{}\textcolor{pink}{Wien I}\oindex{I., Innere Stadt@\textbf{I., Innere Stadt}, \emph{Verwaltungsgebiet}|pw}{}\ledrightnote{\textcolor{pink}{I., Innere Stadt}}\pend{}\pstart{}\textcolor{pink}{Tiefer Graben 17}\oindex{Wien@\textbf{Wien}!I., Innere Stadt@\textbf{I., Innere Stadt}!Tiefer Graben 17@\textbf{Tiefer Graben 17}, \emph{Wohngebäude}|pw}{}\ledrightnote{\textcolor{pink}{Tiefer Graben 17}}\pend{}{\bigskip}
\pstart
           \noindent{}\centering{}{\pb}\textcolor{pink}{\textcolor{gray}{\textbf{Freiburg i. Br.}}}\oindex{Freiburg im Breisgau@\textbf{Freiburg im Breisgau}, \emph{Hauptstadt}|pw}{}\ledrightnote{\textcolor{pink}{Freiburg im Breisgau}}\pend
           
\pstart
           \centering{}\textcolor{pink}{\textcolor{gray}{\textbf{Rathaus}}}\oindex{Rathaus [Freiburg im Breisgau]@\textbf{Rathaus [Freiburg im Breisgau]}, \emph{Verwaltungsgebäude}|pw}{}\ledrightnote{\textcolor{pink}{Rathaus [Freiburg im Breisgau]}}\pend
           \vspace{1em}
\pstart
           \raggedleft{}{\pb}\textcolor{pink}{Freiburg}\oindex{Freiburg im Breisgau@\textbf{Freiburg im Breisgau}, \emph{Hauptstadt}|pw}{}\ledrightnote{\textcolor{pink}{Freiburg im Breisgau}}{ }13. 1. 925\pend
           \vspace{0.5em}
\pstart
           Nun hab ich drei \textcolor{violet}{Vorlesungen}\eventindex{Kultur- und Kongresszentrum Liederhalle@\textbf{Kultur- und Kongresszentrum Liederhalle}!Lesung von Die dreifache Warnung, Das Schicksal des Freiherrn von Leisenbohg, Große Szene, 9.1.1925@Lesung von Die dreifache Warnung, Das Schicksal des Freiherrn von Leisenbohg, Große Szene, 9.1.1925|pwv}\eventindex{Städtische Schauspiele Baden-Baden@\textbf{Städtische Schauspiele Baden-Baden}!Lesung von Das Tagebuch der Redegonda, Die letzten Masken, Weihnachts-Einkäufe, 11.1.1925@Lesung von Das Tagebuch der Redegonda, Die letzten Masken, Weihnachts-Einkäufe, 11.1.1925|pwv}\eventindex{Paulussaal@\textbf{Paulussaal}!Lesung von Das Tagebuch der Redegonda, Die letzten Masken, Große Szene, 12.1.1925@Lesung von Das Tagebuch der Redegonda, Die letzten Masken, Große Szene, 12.1.1925|pwv}{}\ledrightnote{{$\rightarrow$}\emph{\textcolor{violet}{Lesung von Die dreifache Warnung, Das Schicksal des Freiherrn von Leisenbohg, Große Szene, 9.1.1925}}{\newline}{$\rightarrow$}\emph{\textcolor{violet}{Lesung von Das Tagebuch der Redegonda, Die letzten Masken, Weihnachts-Einkäufe, 11.1.1925}}{\newline}{$\rightarrow$}\emph{\textcolor{violet}{Lesung von Das Tagebuch der Redegonda, Die letzten Masken, Große Szene, 12.1.1925}}}{ }\label{K_L04356-1v}\edtext{\textcolor{violet}{\textcolor{pink}{Stuttgart}\oindex{Stuttgart@\textbf{Stuttgart}|pw}{}\ledrightnote{\textcolor{pink}{Stuttgart}}}\eventindex{Kultur- und Kongresszentrum Liederhalle@\textbf{Kultur- und Kongresszentrum Liederhalle}!Lesung von Die dreifache Warnung, Das Schicksal des Freiherrn von Leisenbohg, Große Szene, 9.1.1925@Lesung von Die dreifache Warnung, Das Schicksal des Freiherrn von Leisenbohg, Große Szene, 9.1.1925|pwv}{}\ledrightnote{{$\rightarrow$}\emph{\textcolor{violet}{Lesung von Die dreifache Warnung, Das Schicksal des Freiherrn von Leisenbohg, Große Szene, 9.1.1925}}}}{\lemma{\textnormal{\emph{Stuttgart}}}\Cendnote{\textnormal{Vgl. A. S.: \emph{Kulturveranstaltungen}, 9. 1. 1925.}}}\label{K_L04356-1}, \label{K_L04356-2v}\edtext{\textcolor{violet}{\textcolor{pink}{B.-B.}\oindex{Baden-Baden@\textbf{Baden-Baden}|pw}{}\ledrightnote{\textcolor{pink}{Baden-Baden}}}\eventindex{Städtische Schauspiele Baden-Baden@\textbf{Städtische Schauspiele Baden-Baden}!Lesung von Das Tagebuch der Redegonda, Die letzten Masken, Weihnachts-Einkäufe, 11.1.1925@Lesung von Das Tagebuch der Redegonda, Die letzten Masken, Weihnachts-Einkäufe, 11.1.1925|pwv}{}\ledrightnote{{$\rightarrow$}\emph{\textcolor{violet}{Lesung von Das Tagebuch der Redegonda, Die letzten Masken, Weihnachts-Einkäufe, 11.1.1925}}}}{\lemma{\textnormal{\emph{B.-B.}}}\Cendnote{\textnormal{Vgl. A. S.: \emph{Kulturveranstaltungen}, 11. 1. 1925.}}}\label{K_L04356-2} u \label{K_L04356-3v}\edtext{\textcolor{violet}{\textcolor{pink}{Freiburg}\oindex{Freiburg im Breisgau@\textbf{Freiburg im Breisgau}, \emph{Hauptstadt}|pw}{}\ledrightnote{\textcolor{pink}{Freiburg im Breisgau}}}\eventindex{Paulussaal@\textbf{Paulussaal}!Lesung von Das Tagebuch der Redegonda, Die letzten Masken, Große Szene, 12.1.1925@Lesung von Das Tagebuch der Redegonda, Die letzten Masken, Große Szene, 12.1.1925|pwv}{}\ledrightnote{{$\rightarrow$}\emph{\textcolor{violet}{Lesung von Das Tagebuch der Redegonda, Die letzten Masken, Große Szene, 12.1.1925}}}}{\lemma{\textnormal{\emph{Freiburg}}}\Cendnote{\textnormal{Vgl. A. S.: \emph{Kulturveranstaltungen}, 12. 1. 1925.}}}\label{K_L04356-3} sehr gut hinter mir – heute  Abend
                  begi{\geminationn} ich, mit \label{K_L04356-4v}\edtext{\textcolor{violet}{\textcolor{pink}{Basel}\oindex{Basel@\textbf{Basel}|pw}{}\ledrightnote{\textcolor{pink}{Basel}}}\eventindex{Stadtcasino Basel@\textbf{Stadtcasino Basel}!Lesung von Das Tagebuch der Redegonda, Die letzten Masken, Große Szene, 13.1.1925@Lesung von Das Tagebuch der Redegonda, Die letzten Masken, Große Szene, 13.1.1925|pwv}{}\ledrightnote{{$\rightarrow$}\emph{\textcolor{violet}{Lesung von Das Tagebuch der Redegonda, Die letzten Masken, Große Szene, 13.1.1925}}}}{\lemma{\textnormal{\emph{Basel}}}\Cendnote{\textnormal{Vgl. A. S.: \emph{Kulturveranstaltungen}, 13. 1. 1925.}}}\label{K_L04356-4}, die \textcolor{pink}{Schweiz}\oindex{Schweiz@\textbf{Schweiz}|pw}{}\ledrightnote{\textcolor{pink}{Schweiz}} zu beglücken.
               Befinden vortrefflich. Lassen Sie mich vielleicht nach \textcolor{pink}{Zürich}\oindex{Zürich@\textbf{Zürich}|pw}{}\ledrightnote{\textcolor{pink}{Zürich}}{ }\textcolor{pink}{Hotel Gotthard}\oindex{Hotel St. Gotthard Zürich@\textbf{Hotel St. Gotthard Zürich}, \emph{Hotel}|pw}{}\ledrightnote{\textcolor{pink}{Hotel St. Gotthard Zürich}} mit einer
                  Zeil\textcolor{gray}{e} wissen, wies bei Ihnen geht.\pend
           
\pstart
           Alles{\\[\baselineskip]}herzliche. {[}Ihr Arthur{]}\pend
           \leftskip=0em{}\selectlanguage{ngerman}\endnumbering\briefempfaengerindex{Schwarzkopf, Gustav@\textsc{Schwarzkopf, Gustav}!zzzSchnitzler, Arthur@\emph{von Arthur Schnitzler}!1925-01-131@{13. 1. 1925}|)be}\mylabel{L04356h}
\begin{anhang}
\end{anhang}\newcommand{\dateiname}{L04356}\newcommand{\titel}{Arthur Schnitzler an Gustav Schwarzkopf, 13. 1. 1925}\newcommand{\editorInnen}{Herausgegeben von Jahnke, SelmaMüller, Martin Anton}\input{../tex-inputs/latex-korrekturansicht-abspann}
